\documentclass[11pt,a4paper]{article}

\usepackage[utf8]{inputenc}
\usepackage{amsmath,amssymb,mathtools}
\usepackage{geometry}
\usepackage{microtype}

\geometry{margin=1in}

\title{Analytical Pixel Integration of Gaussian Speckle Blobs}
\author{Renderer Convergence Conjecture}
\date{\today}

\begin{document}

\maketitle

\section{Purpose}

This note derives the pixel integral for the untruncated Gaussian blobs used
in Experiment 2.  It distinguishes the rigid and general affine cases, and
records the consequences of the current radial cutoff and additive clamp.

\section{Reference Speckle Field}

Let $\mathbf{r}=(x,y)^T$ be reference coordinates, and let blob $j$ have
centre $\mathbf{c}_j$ and common standard deviation $\sigma$.  Its
untruncated coverage is
\begin{equation}
 g_j(\mathbf{r}) =
 \exp\left(-\frac{\lVert\mathbf{r}-\mathbf{c}_j\rVert^2}
 {2\sigma^2}\right).
\end{equation}
Ignoring overlap saturation for the moment, the total coverage and intensity
are
\begin{align}
 G(\mathbf{r}) &= \sum_j g_j(\mathbf{r}), \\
 I(\mathbf{r}) &= I_0 + \gamma\left[1-2G(\mathbf{r})\right].
\end{align}
Thus a unit coverage is black, $I_0-\gamma$, and zero coverage is white,
$I_0+\gamma$.  The implemented renderer instead uses
$G \leftarrow \min(G,1)$; this is discussed in Section~\ref{sec:limits}.

\section{Pixel Mapping}

Let $\mathbf{q}=(q_x,q_y)^T\in[-h,h]^2$ be coordinates relative to the
centre of a target pixel, where $2h$ is the target-pixel width.  On a rigid
or affine patch, the inverse map from target to reference coordinates is
\begin{equation}
 \mathbf{r}(\mathbf{q}) = \mathbf{r}_0 + \mathbf{A}\mathbf{q}.
 \label{eq:affine-map}
\end{equation}
Here $\mathbf{A}$ is the local inverse deformation gradient and
$\mathbf{r}_0$ is the reference coordinate corresponding to the target
pixel centre.  The pixel average of blob $j$ is
\begin{equation}
 \bar{g}_j = \frac{1}{4h^2}
 \int_{-h}^{h}\int_{-h}^{h}
 g_j\left(\mathbf{r}_0+\mathbf{A}\mathbf{q}\right)
 \,dq_y\,dq_x.
 \label{eq:pixel-average}
\end{equation}

Define $\mathbf{d}_j=\mathbf{r}_0-\mathbf{c}_j$.  If $\mathbf{A}$ is
nonsingular, define
\begin{equation}
 \boldsymbol{\mu}_j=-\mathbf{A}^{-1}\mathbf{d}_j,
 \qquad
 \mathbf{B}=\mathbf{A}^T\mathbf{A},
 \qquad
 \boldsymbol{\Sigma}=\sigma^2\mathbf{B}^{-1}.
\end{equation}
The exponent in Equation~\eqref{eq:pixel-average} becomes
\begin{equation}
 g_j\left(\mathbf{r}(\mathbf{q})\right)=
 \exp\left[-\frac{1}{2}
 (\mathbf{q}-\boldsymbol{\mu}_j)^T
 \mathbf{B}(\mathbf{q}-\boldsymbol{\mu}_j)/\sigma^2\right].
 \label{eq:target-gaussian}
\end{equation}

\section{General Affine Result}

Let $\Phi_2(\mathbf{z};\boldsymbol{\mu},\boldsymbol{\Sigma})$ denote the
bivariate normal CDF through upper limits $\mathbf{z}$.  For the rectangle
$R=[-h,h]^2$, let $P_R$ be the usual inclusion--exclusion of its four CDF
corners.  Then
\begin{equation}
 \int_R g_j\left(\mathbf{r}(\mathbf{q})\right)\,d\mathbf{q}
 = \frac{2\pi\sigma^2}{\left|\det\mathbf{A}\right|}P_R.
 \label{eq:affine-cdf}
\end{equation}
Equation~\eqref{eq:affine-cdf} is exact for any nonsingular affine map.
For a general shear, it requires a correlated bivariate normal CDF, rather
than only the scalar error function.

\section{Error-Function Form}

If $\mathbf{B}$ is diagonal, write
\begin{equation}
 \mathbf{B}=\operatorname{diag}(b_x,b_y).
\end{equation}
The target-coordinate Gaussian is then separable.  Define
\begin{align}
 E_x &= \operatorname{erf}\left(
 \frac{\sqrt{b_x}(h-\mu_{j,x})}{\sqrt{2}\sigma}
 \right)-\operatorname{erf}\left(
 \frac{\sqrt{b_x}(-h-\mu_{j,x})}{\sqrt{2}\sigma}
 \right), \\
 E_y &= \operatorname{erf}\left(
 \frac{\sqrt{b_y}(h-\mu_{j,y})}{\sqrt{2}\sigma}
 \right)-\operatorname{erf}\left(
 \frac{\sqrt{b_y}(-h-\mu_{j,y})}{\sqrt{2}\sigma}
 \right).
\end{align}
The exact blob integral is
\begin{equation}
 \int_R g_j\left(\mathbf{r}(\mathbf{q})\right)\,d\mathbf{q}
 = \frac{\pi\sigma^2}{2\sqrt{b_xb_y}} E_xE_y.
 \label{eq:erf-result}
\end{equation}

For a rigid map, $\mathbf{A}$ is a rotation and therefore
$\mathbf{B}=\mathbf{I}$, so Equation~\eqref{eq:erf-result} applies
directly.  Translation-only rigid motion is the special case
$\mathbf{A}=\mathbf{I}$.  Axis-aligned affine scaling also remains
separable.  A general affine map with shear or non-orthogonal mapped axes
does not: use Equation~\eqref{eq:affine-cdf} instead.

\section{Intensity Integral}

If the blobs do not saturate when added, linearity gives
\begin{equation}
 \bar{I}=I_0+\gamma\left[1-2\sum_j\bar{g}_j\right].
\end{equation}
Only nearby blobs need be considered in a numerical implementation.  The
untruncated expression has nonzero tails, but they decay exponentially and
may be culled by a conservative tolerance without changing the formula.

\section{Cutoff and Clamp Limitations}
\label{sec:limits}

The present pattern truncates each blob radially at $k\sigma$:
\begin{equation}
 g_j^{(k)}(\mathbf{r})=g_j(\mathbf{r})
 \mathbf{1}\!\left(\lVert\mathbf{r}-\mathbf{c}_j\rVert\leq k\sigma\right).
\end{equation}
Under Equation~\eqref{eq:affine-map}, this adds an elliptical domain to the
pixel rectangle.  The integral is consequently over the intersection of a
rectangle and an ellipse, and is not represented by the simple products of
error functions above.  The fraction of a complete untruncated Gaussian's
integral discarded outside the radius is
\begin{equation}
 \exp\left(-k^2/2\right).
\end{equation}
For $k=4$, this is approximately $3.35\times10^{-4}$.  Removing the hard
cutoff is therefore the cleanest route to an analytic reference.

Likewise, the implemented $\min(\sum_j g_j,1)$ is nonlinear.  Where blobs
overlap enough to saturate, the pixel must be partitioned by the saturation
contour, so the linear sum of Equation~\eqref{eq:affine-cdf} is no longer
exact.  The analytic path is exact when saturation is absent; otherwise a
numerical correction or a different non-saturating composition rule is
required.

\section{Implementation Consequence}

For a first analytic Gaussian renderer, use untruncated blobs and the linear
sum.  Use Equation~\eqref{eq:erf-result} for rigid motion and diagonal
affine maps.  For a fully general affine deformation, evaluate the four
terms of Equation~\eqref{eq:affine-cdf} with a robust bivariate-normal CDF.
Keep the current quadrature renderer as the reference for cutoff or clamp
behaviour until their domain intersections are handled explicitly.

\end{document}
